\documentclass[
  fontsize=14,
  DIV=15,
]{scrreprt}

\usepackage[no-math]{fontspec}

\usepackage{polyglossia}
\setdefaultlanguage[indentfirst=true,spelling=modern]{russian}
\setotherlanguage{english}

\usepackage{amsmath}
\usepackage{unicode-math}

\unimathsetup{
  mathbf=sym,
  mathrm=sym,
}

\setmainfont{CMU Serif}
\setromanfont{CMU Serif}
\setsansfont{CMU Sans Serif}
\setmonofont{CMU Typewriter Text}
\setmathfont{Latin Modern Math}

\usepackage{hyperref}
\hypersetup{
  colorlinks,
  citecolor=black,
  filecolor=black,
  linkcolor=black,
  urlcolor=black,
  unicode=true,
  pdftitle={ВКР Коннова — Модели конкуренции и PINN},
  pdfauthor={Коннова Т.}
}

\usepackage{minted}
\usemintedstyle{friendly}
\setminted{
  autogobble=true,
  breaklines=true,
  breakanywhere=true,
  obeytabs=false,
  tabsize=2,
  linenos=true,
}
\usepackage{verbatim}

\usepackage{graphicx}
\graphicspath{{img/}}

\usepackage[autostyle]{csquotes}
\usepackage{url}

\usepackage[%
  backend=biber,
  bibstyle=gost-numeric,
  defernumbers=true,
  movenames=false,
  maxbibnames=99,
  sorting=nyt
]{biblatex}
\bibliography{lit}

\usepackage{booktabs}
\usepackage{caption}
\usepackage{subcaption}


\begin{document}

\begin{titlepage}
  \centering
  \vspace*{1cm}
  {\large Российский университет дружбы народов имени Патриса Лумумбы}\\[0.5cm]
  {\large Факультет физико-математических и естественных наук}\\[0.5cm]
  {\large Кафедра теории вероятности и кибербезопасности}\\[2cm]

  {\LARGE\bfseries Сравнительное исследование численных и нейросетевых методов моделирования динамики конкурирующих популяций}\\[1.5cm]

  {\large Выпускная квалификационная работа магистра}\\[2cm]

  \begin{flushright}
    {\large Студент: Коннова Т.}\\[0.3cm]
    {\large Научный руководитель: Штепа К. А.}\\
  \end{flushright}

  \vfill
  {\large Москва --- 2027}
\end{titlepage}

\tableofcontents
\newpage

\chapter*{Введение}
\addcontentsline{toc}{chapter}{Введение}

% TODO: Актуальность, цель, задачи, объект и предмет исследования

\chapter{Математическое моделирование популяционной динамики}
\label{ch:models}

\section{Модель Лотки-Вольтерры}
\label{sec:lv-model}

% Предпосылки построения, отличие конкуренции от модели «хищник-жертва»

\section{Классическая модель конкуренции}
\label{sec:competition-ode}

% Система ОДУ, переменные, коэффициенты

\section{Режимы конкуренции}
\label{sec:regimes}

% Сосуществование, вытеснение, бистабильность

\section{Диффузионная модель конкуренции}
\label{sec:diffusion-pde}

% Математическая постановка PDE, начальные и граничные условия


\chapter{Классические численные методы и PINN}
\label{ch:methods}

\section{Численные методы решения ОДУ}
\label{sec:rk-methods}

% Методы Рунге-Кутты, оценка погрешности

\section{Численные методы решения УЧП}
\label{sec:pde-methods}

% Метод конечных разностей, метод прямых

\section{Физически-информированные нейронные сети (PINN)}
\label{sec:pinn}

% Принципы PINN, архитектура, функция потерь

\section{Применение PINN для решения ОДУ}
\label{sec:pinn-ode}

% NNODE, формирование функции потерь

\section{Применение PINN для решения УЧП}
\label{sec:pinn-pde}

% PhysicsInformedNN

\section{Методика сравнительного исследования}
\label{sec:methodology}

% Критерии оценки, метрики ошибок


\chapter{Программная реализация и вычислительные эксперименты}
\label{ch:experiments}

\section{Выбор инструментальных средств}
\label{sec:tools}

% Julia, DifferentialEquations.jl, NeuralPDE.jl, Lux.jl

\section{Численное решение классической модели конкуренции}
\label{sec:exp-ode-numeric}

\section{Решение классической модели с использованием PINN}
\label{sec:exp-ode-pinn}

\section{Численное решение диффузионной модели}
\label{sec:exp-pde-numeric}

\section{Решение диффузионной модели с использованием PINN}
\label{sec:exp-pde-pinn}

\section{Исследование различных режимов конкуренции}
\label{sec:exp-regimes}

\section{Сравнительный анализ точности и вычислительных затрат}
\label{sec:exp-comparison}

\section{Анализ полученных результатов}
\label{sec:exp-analysis}


\chapter*{Заключение}
\addcontentsline{toc}{chapter}{Заключение}

% Основные результаты, выводы, направления дальнейшего развития

\printbibliography[title={Список литературы}]
\addcontentsline{toc}{chapter}{Список литературы}

\end{document}
